\chapter{Réalisation technique et Implémentation}

\section{Mise en œuvre de l'architecture 3-tiers}

Ce chapitre détaille la phase de réalisation de l'application \textbf{My Smart Budget}. Il présente la traduction des choix architecturaux (vus au chapitre précédent) en code fonctionnel, en respectant la séparation stricte des trois couches : Présentation (Next.js), Métier (Node.js/Express) et Données (PostgreSQL/MongoDB).

\subsection{Vue d'ensemble de la stack}

L'implémentation repose sur une stack moderne et robuste, choisie pour sa flexibilité et sa performance :

\begin{itemize}
    \item \textbf{Tier 1 (Client) :} Next.js 15 + React 19 + TypeScript.
    \item \textbf{Tier 2 (Serveur) :} Node.js + Express.js (API REST).
    \item \textbf{Tier 3 (Données) :} PostgreSQL (Transactionnel) + MongoDB (Logs/Analytics).
\end{itemize}

\begin{figure}[H]
    \centering
    \safefig{chapters/assets/images/architecture_3_tiers.png}{0.90\textwidth}{architecture_3_tiers.png}
    \caption{Flux de données implémenté dans l'architecture 3-tiers}
    \label{fig:arch_impl}
\end{figure}

\section{Couche Présentation (Frontend)}

L'interface utilisateur a été développée avec une approche par composants, garantissant la réutilisabilité et la maintenabilité du code.

\subsection{Structure et Composants}

L'arborescence du projet Frontend reflète le découpage fonctionnel de l'application :

\begin{exemple}
\textbf{Organisation du code source (src/) :}
\begin{verbatim}
src/
+-- components/
|   +-- common/            # Composants UI génériques (Input, Button...)
|   +-- dashboard/         # Widgets (Charts, KPIs)
|   +-- transactions/      # Listes et Formulaires CRUD
|   +-- budgets/           # Gestion des enveloppes
+-- hooks/                 # Logique réutilisable (Custom Hooks)
|   +-- useTransactions.ts
|   +-- useBudgetAlerts.ts
+-- store/                 # Gestion d'état (React Context + useState)
+-- services/              # Appels API (fetch natif / Axios)
\end{verbatim}
\end{exemple}

\subsection{Implémentation des formulaires complexes}

La gestion des transactions nécessite une validation rigoureuse côté client pour garantir l'intégrité des données avant l'envoi au serveur. L'exemple ci-dessous montre le composant principal de saisie, avec validation locale et avertissement en cas de dépassement de budget :

\begin{lstlisting}[language=JavaScript, caption={Composant TransactionForm avec validation et UX}]
const TransactionForm = ({ onSubmit, categories, budget }) => {
  const [formData, setFormData] = useState({
    amount: '',
    category: '',
    date: new Date()
  });
  
  const handleSubmit = async (e) => {
    e.preventDefault();
    
    // 1. Validation locale
    if (parseFloat(formData.amount) <= 0) {
      setError('Le montant doit être positif');
      return;
    }
    
    // 2. Vérification proactive du budget (UX)
    if (budget && parseFloat(formData.amount) > budget.remaining) {
      const confirm = await showConfirmDialog(
        'Attention : Dépassement de budget',
        `Cette dépense va excéder votre plafond de ${
          parseFloat(formData.amount) - budget.remaining
        }€. Continuer ?`
      );
      if (!confirm) return;
    }
    
    // 3. Soumission
    await onSubmit(formData);
  };

  return (
    <form onSubmit={handleSubmit} aria-label="Ajouter une transaction">
      {/* Champs natifs stylisés Tailwind CSS */}
      <input
        type="number"
        placeholder="Montant (€)"
        value={formData.amount}
        onChange={handleChange}
        className="border rounded px-3 py-2 w-full"
      />
      {/* ... Selecteur de catégorie ... */}
      <button type="submit" className="bg-blue-600 text-white px-4 py-2 rounded">Enregistrer</button>
    </form>
  );
};
\end{lstlisting}

Ce composant garantit que l'utilisateur est averti avant tout dépassement de budget, sans bloquer l'action si elle est confirmée.

\section{Couche Logique Métier (Backend)}

Le serveur API, développé avec Express.js, centralise la logique métier. Il agit comme chef d'orchestre entre le client et les bases de données.

\subsection{Contrôleurs (Controllers)}

Les contrôleurs gèrent les requêtes HTTP, valident les entrées et délèguent le traitement aux services.

\begin{exemple}
\textbf{TransactionController.js — Traitement d'une création :}
\begin{lstlisting}[language=JavaScript]
class TransactionController {
  async createTransaction(req, res) {
    try {
      const { amount, category, date, type } = req.body;
      const userId = req.user.id; // Extrait du JWT
      
      // 1. Appel au Service Métier
      const transaction = await this.transactionService.create({
        userId, amount, category, date, type
      });
      
      // 2. Déclenchement asynchrone des vérifications
      // On n'attend pas la fin pour répondre au client (Performance)
      this.budgetService.checkBudgetAlert(userId, category, new Date())
        .catch(err => console.error('Alert Error:', err));
      
      res.status(201).json(transaction);
    } catch (error) {
      res.status(400).json({ error: error.message });
    }
  }
}
\end{lstlisting}
\end{exemple}

\subsection{Services Métier}

C'est ici que réside l'intelligence de l'application (calculs, règles de gestion). L'exemple suivant montre comment le service surveille automatiquement la consommation d'un budget après chaque transaction :

\begin{exemple}
\textbf{BudgetService.js — Logique d'alerte :}
\begin{lstlisting}[language=JavaScript]
class BudgetService {
  async checkBudgetAlert(userId, categoryId, date) {
    // Récupération du budget actif
    const budget = await this.budgetRepo.findActive(userId, categoryId, date);
    if (!budget) return;
    
    // Calcul de la consommation en temps réel
    const spent = await this.transactionRepo.sumByCategory(
      userId, categoryId, budget.startDate, budget.endDate
    );
    
    const percentage = (spent / budget.amount) * 100;
    
    // Règle métier : Alerte si seuil dépassé
    if (percentage >= budget.alertThreshold) {
      await this.notificationService.send({
        userId,
        type: 'BUDGET_WARNING',
        message: `Attention, vous avez consommé ${percentage.toFixed(1)}% de votre budget ${budget.name}.`
      });
    }
  }
}
\end{lstlisting}
\end{exemple}

Concrètement, chaque ajout de dépense déclenche silencieusement cette vérification : l'utilisateur reçoit une notification sans même avoir à consulter son tableau de bord.

\section{Couche Données (Data Layer)}

L'application utilise une stratégie de persistance hybride (Polyglot Persistence) pour optimiser les performances et la traçabilité.

\subsection{Modélisation Relationnelle (PostgreSQL)}

Les données critiques (Transactions, Budgets) nécessitent une intégrité référentielle forte (ACID).

\begin{lstlisting}[language=SQL, caption={Schéma SQL des transactions}]
CREATE TABLE transactions (
  id UUID PRIMARY KEY DEFAULT gen_random_uuid(),
  user_id UUID NOT NULL REFERENCES users(id),
  amount DECIMAL(10,2) NOT NULL, -- DECIMAL pour la précision monétaire
  type VARCHAR(10) CHECK (type IN ('income', 'expense')),
  category_id UUID REFERENCES categories(id),
  date TIMESTAMP NOT NULL,
  created_at TIMESTAMP DEFAULT CURRENT_TIMESTAMP
);

-- Index pour optimiser le Dashboard (filtrage par date)
CREATE INDEX idx_transactions_user_date ON transactions(user_id, date DESC);
\end{lstlisting}

\subsection{Modélisation Documentaire (MongoDB)}

MongoDB est utilisé pour stocker les logs d'audit et les rapports analytiques, dont la structure peut évoluer.

\begin{lstlisting}[language=JavaScript, caption={Schéma Mongoose pour les rapports}]
const monthlyReportSchema = {
  userId: ObjectId,
  period: { month: Number, year: Number },
  metrics: {
    totalExpenses: Number,
    savingsRate: Number,
    categoryBreakdown: [{ 
      categoryName: String, 
      amount: Number, 
      percentage: Number 
    }]
  },
  generatedAt: Date // TTL index possible pour suppression auto
};
\end{lstlisting}

\section{Sécurité et Conformité}

La sécurité a été intégrée dès la conception ("Security by Design") pour protéger les données financières.

\subsection{Validation et Assainissement}

Toutes les données entrantes passent par une double barrière de validation.

\begin{lstlisting}[language=JavaScript, caption={Middleware de validation Express}]
const validateTransaction = [
  // Validation du type et de la plage de valeur
  body('amount')
    .isFloat({ min: 0.01, max: 1000000 })
    .withMessage('Montant invalide'),
  
  // Assainissement (Sanitization) pour éviter les XSS
  body('description')
    .trim()
    .escape(),
    
  (req, res, next) => {
    const errors = validationResult(req);
    if (!errors.isEmpty()) return res.status(400).json({ errors: errors.array() });
    next();
  }
];
\end{lstlisting}

\subsection{Chiffrement des données sensibles}

Les informations critiques (comme les tokens bancaires) sont chiffrées en base via AES-256-GCM.

\begin{lstlisting}[language=JavaScript]
encryptData(text) {
  const iv = crypto.randomBytes(16);
  const cipher = crypto.createCipheriv('aes-256-gcm', process.env.KEY, iv);
  let encrypted = cipher.update(text, 'utf8', 'hex');
  encrypted += cipher.final('hex');
  return { encrypted, iv: iv.toString('hex'), tag: cipher.getAuthTag() };
}
\end{lstlisting}

\section{Bilan de la réalisation}

L'implémentation respecte l'architecture définie lors de la conception. L'utilisation de \textbf{Services} découplés des \textbf{Contrôleurs} facilite les tests unitaires, tandis que la séparation Frontend/Backend via une API REST permet d'envisager sereinement le développement futur d'une application mobile native utilisant les mêmes endpoints.

En résumé, cette couche d'implémentation traduit fidèlement les choix de conception en code fonctionnel, testé et sécurisé. Ce que le jury doit retenir : chaque couche a une responsabilité claire, et aucune logique métier ne s'est glissée dans les contrôleurs ou les composants Next.js.